\scanpage{23}
on trouve une \deleted{\unsure{plongée|plongement}} immersion canonique
\[
\Omega\hookrightarrow E \tag*{(6)}
\]
également compatible à tt chgt. de base. Soit
\deleted{\unsure{une application}}\deleted{\unsure{le couple plongement et}} $F$ un fibré affine sous
un fibré vectoriel $\check{\mathbb V}(V)$, \deleted{et considérons un}\inserted{défini par une}
\deleted{\unsure{$S$-morphisme}}\inserted{structure} d'extension
\[
0\mathbin{-}V\longrightarrow\check C\longrightarrow\underline{O}_S\longrightarrow0 \tag*{(7)}
\]
\unsure{ou|\textemdash}
\[
0\mathbin{-}\underline{O}_S\longrightarrow C^{\deleted{\unsure{\vee}}}\mathbin{-}\check V\mathbin{-}0 \tag*{\deleted{(8)}}
\]
avec
\[
F=\operatorname{Spec}(\underline{\mathrm{Sym}}\, C^{\deleted{\unsure{\vee}}}/(e_C-1))
\]
où $e_C$ est l'image \inserted{dans $C$} de la section 1 de $\underline{O}_S$.

Mais les hom
\[
\Omega\longrightarrow F
\]
correspondent aux hom de $\underline{O}_S$-algèbres
\[
\underline{\mathrm{Sym}}(C)/(e_C-1)\longrightarrow B
\]
ou encore aux hom de $\underline{O}_S$-Modules
\[
C\xrightarrow{\varphi}B
\]
tels que
\[
\varphi(e_C)=e_B
\]
ou, par transposition, aux hom d'extensions de
(3) dans (7), i.e. aux \uline{hom} \uline{de fibrés affines}
$E\mathbin{-}F$. Donc on trouve un isom canonique
\[
\operatorname{Hom}_S(\Omega,F)\simeq\operatorname{Hom\,aff}_S(E,F) \tag*{(9)}
\]
(ce qui justifie le nom d'\uline{enveloppe affine de $\Omega$}
donné à $E$), d'où par chgt de base
\[
\underline{\mathrm{Hom}}_S(\Omega,F)\simeq\underline{\mathrm{Hom}}\,\mathrm{aff}_S(E,F).\tag*{\deleted{\unsure{ou}}(10)}
\]

% archive sheet number: 22

